\subsection{Information-channel test with CHFS}%
%\label{subsec:chfs_info_channel}

The aggregate horse race in Section~\ref{subsec:horse_race} operates
with 32 quarterly observations and cannot cleanly separate diagnostic
overreaction from information rigidity. Both channels predict a
negative revision-error coefficient, and multicollinearity between the
revision and lagged-error terms makes the joint specification
indeterminate at this sample size. The CHFS 2011 wave provides a
different discriminating test that does not require time-series
variation: it exploits cross-sectional differences in how households
acquire price information.

\paragraph{Discriminating logic.}
The diagnostic model predicts that overreaction intensifies when
households receive information through high-salience channels.
Television news, newspaper headlines, and internet portals frame price
movements as narratives with emotional valence, exactly the type of
signal that triggers representativeness-based overweighting
\citep{Bordalo2018}. Households who rely instead on family, friends, or
SMS alerts receive more diffuse signals with weaker narrative framing.
The prediction is that media-intensive households overreact more.

Information rigidity makes the opposite prediction. If the primary
friction is that households update too slowly because they pay
insufficient attention to public signals, then households with better
access to media should update more completely and make smaller forecast
errors. Media exposure reduces the information gap. The two models
therefore generate opposing signs on the interaction between
expectations and media intensity.

\paragraph{Specification.}
I estimate three specifications with increasing stringency. The
baseline tests whether media exposure predicts the level of inflation
expectations:
\begin{equation}%\label{eq:chfs_info_channel}
  Exp_{ip} = \alpha + \beta_1 HighMedia_i
    + \gamma X_i + \delta_p + \varepsilon_{ip},
\end{equation}
where $Exp_{ip}$ is the inflation expectation for household $i$ in
province $p$ on the five-point A4008 scale (recoded so higher values
indicate more inflationary beliefs). $HighMedia_i$ is an indicator
equal to one if the household reports above-median use of TV, internet,
or newspapers as information sources in A4001, and zero otherwise.
$X_i$ includes log household income and an urban/rural indicator.
$\delta_p$ denotes province fixed effects.

The key coefficient is $\beta_1$. In 2011, when actual CPI inflation
was elevated at approximately 5.4\% but moderating, the rational
expectation maps to ``rise a little'' (ordinal value 4 on the recoded
scale). Diagnostic expectations predict $\beta_1 > 0$: households
exposed to high-salience media channels form more inflationary beliefs
that overshoot the rational benchmark. Information rigidity predicts
$\beta_1 \leq 0$: media-exposed households receive more information
and therefore hold more moderate, accurate beliefs.

A second specification replaces the level with expectation extremity
(absolute deviation from the province mean), testing whether media
exposure amplifies the dispersion of beliefs. A third uses a binary
indicator for the most extreme category (``rise a lot'') as the
dependent variable.

\paragraph{Results.}
Table~\ref{tab:chfs_mechanism_info_channel} reports the estimates.
Across all three dependent variables, the high-media coefficient is
small and statistically indistinguishable from zero. Column~1 shows
$\hat{\beta}_1 = -0.005$ (SE $= 0.022$) for the expectation level.
Column~2 shows $-0.001$ (SE $= 0.015$) for extremity. Column~3 shows
$-0.004$ (SE $= 0.012$) for the over-prediction binary. None
approaches conventional significance. Column~4 reports the risk
attitude coefficient, which is also null ($-0.010$, SE $= 0.008$).
Column~5 enters both media and risk attitude jointly; neither is
significant.

\begin{table}[htbp]
\centering
\begin{threeparttable}
\caption{CHFS 2011: Information Channels and Inflation Expectations}
\label{tab:chfs_mechanism_info_channel}
{\small
\begin{tabular}{lccccc}
\toprule
 & (1) & (2) & (3) & (4) & (5) \\
 & Infl.\ exp. & Extremity & Over-predict & Risk placebo & Horse race \\
\midrule
\multicolumn{6}{l}{\textit{Panel A: Information Channel}} \\
\addlinespace
High media (TV/internet/newspaper) & $-0.005$ & $-0.001$ & $-0.004$ & & $-0.002$ \\
 & $(0.022)$ & $(0.015)$ & $(0.012)$ & & $(0.022)$ \\
 & $[0.820]$ & $[0.947]$ & $[0.739]$ & & $[0.928]$ \\
\addlinespace
\multicolumn{6}{l}{\textit{Panel B: Risk Preference Placebo}} \\
\addlinespace
Risk attitude & & & & $-0.010$ & $-0.010$ \\
 & & & & $(0.008)$ & $(0.008)$ \\
 & & & & $[0.211]$ & $[0.211]$ \\
\addlinespace
Observations & 8,281 & 8,281 & 8,281 & 8,173 & 8,173 \\
Province FE & Yes & Yes & Yes & Yes & Yes \\
Controls & Yes & Yes & Yes & Yes & Yes \\
\bottomrule
\end{tabular}
}
\begin{tablenotes}[flushleft]
\footnotesize
\item \textit{Notes:} Columns (1), (4), and (5) use inflation expectation (1--5 scale, higher = more inflationary) as the dependent variable. Column (2) uses absolute deviation from province mean expectation. Column (3) uses an indicator for expecting prices to rise a lot. Controls include log income and an urban indicator. Heteroskedasticity-robust standard errors are in parentheses; p-values are in brackets. $^{*}p<0.10$; $^{**}p<0.05$; $^{***}p<0.01$.
\end{tablenotes}
\end{threeparttable}
\end{table}


\begin{figure}[!htbp]
  \centering
  \includegraphics[width=0.75\textwidth]{chfs_info_channel_gradient}
  \caption{Mean inflation expectations by media-exposure group,
    CHFS 2011. Higher values indicate more inflationary beliefs
    (5-point scale). The unconditional gap between high-media and
    low-media households is small and largely absorbed by province
    fixed effects.}
  %\label{fig:chfs_info_channel}
\end{figure}

The null result has a clear interpretation. Province fixed effects
absorb the dominant source of variation in inflation expectations:
local price conditions. Within provinces, media exposure does not
predict belief intensity. Figure~\ref{fig:chfs_info_channel} confirms
this visually: the unconditional gap in mean expectations between
high-media and low-media households (4.14 versus 4.08) is small and
largely explained by geographic sorting.

\paragraph{Why the cross-sectional test is null and the panel test is not.}
The CFPS panel result ($\hat{\beta} \approx -0.55$, $p < 0.001$)
identifies overreaction through within-household \textit{temporal}
variation: the same household revises expectations and then
overshoots. The CHFS cross-sectional test asks a different question:
whether households with different \textit{persistent} information
environments hold systematically different beliefs. The null answer
is consistent with the diagnostic model if overreaction is triggered
by salient \textit{events} rather than by the household's general
media diet. A household that watches TV every day does not overreact
more than one that does not, but a household that revises its
expectation after a salient food-price spike overshoots regardless of
its information channel. This distinction aligns with the
representativeness mechanism in \citet{Bordalo2018}: the bias is
event-driven, not channel-driven.

The CHFS test does rule out one specific version of the diagnostic
story: that persistent differences in media access generate
persistent cross-sectional differences in overreaction intensity.
It does not rule out the event-driven version that the CFPS panel
identifies. The two results are complementary, not contradictory.

\begin{table}[htbp]
\centering
\begin{threeparttable}
\caption{CHFS 2011: Risk Attitude Placebo Test}
\label{tab:chfs_risk_placebo}
{\small
\begin{tabular}{lc}
\toprule
 & (1) \\
 & Infl.\ exp. \\
\midrule
\multicolumn{2}{l}{\textit{Panel A: Placebo Outcome}} \\
\addlinespace
Risk attitude & $-0.010$ \\
 & $(0.008)$ \\
 & $[0.211]$ \\
\addlinespace
Observations & 8,173 \\
Province FE & Yes \\
Controls & Yes \\
\bottomrule
\end{tabular}
}
\begin{tablenotes}[flushleft]
\footnotesize
\item \textit{Notes:} The dependent variable is inflation expectation (1--5, higher = more inflationary). Risk attitude should not predict expectations if the diagnostic channel operates through salience rather than risk preferences. Controls include log income and an urban indicator. Heteroskedasticity-robust standard errors are in parentheses; p-values are in brackets. $^{*}p<0.10$; $^{**}p<0.05$; $^{***}p<0.01$.
\end{tablenotes}
\end{threeparttable}
\end{table}


\paragraph{Interpretation and caveats.}
The CHFS information-channel test is cross-sectional. With a single
wave, it captures persistent household differences in information
environment but cannot isolate the temporal salience shocks that drive
the CFPS panel result. The 2011 wave was fielded during a period of
elevated but moderating inflation, when virtually all households
expected prices to rise. This ceiling effect compresses the
distribution and reduces the scope for media channels to generate
differential beliefs. A richer test would exploit within-household
variation in media exposure to salient events at higher frequency,
which the current data do not support.
